\documentclass[a4paper,12pt]{article}
\usepackage[fontset=none]{ctex}
\usepackage{geometry}
\usepackage{graphicx}
\usepackage{amsmath}
\usepackage{enumitem}
\usepackage{setspace}
\usepackage{fontspec}

% 字体设置（与 docx 模板一致）
\setCJKmainfont{Songti SC}          % 宋体 - 正文默认
\setCJKsansfont{Heiti SC}            % 黑体 - 一级标题
\setCJKmonofont[AutoFakeBold]{STFangsong}          % 仿宋 - 注意事项
\setmainfont{Times New Roman}        % 英文默认

% 页面设置
\geometry{left=2.5cm, right=2.5cm, top=2.5cm, bottom=2.5cm}

% 去掉页码
\pagestyle{empty}

\begin{document}

% ========== 封面 ==========
\begin{center}
    \vspace*{2cm}
    {\fontsize{36pt}{43pt}\selectfont \textbf{物理实验预习报告}}
    \vspace{3cm}
\end{center}

\begin{center}
    \fontsize{16pt}{24pt}\selectfont
    \textbf{实验名称：铁磁材料的磁滞回线和基本磁化曲线\quad} \\[0.5em]
    \textbf{指导教师：\_\_\_\_\_\_\_\_\_\_\_\_\_\_\_\_\_\_\_\_\_\_\_\_\_}
\end{center}

\vspace{2cm}

\begin{center}
    \fontsize{14pt}{28pt}\selectfont
    \textbf{周次：\_\_\_\_\_\_\_\_\_\_\_\_\_\_\_\_\_\_\_\_\_\_\_\_\_\_\_\_\_\_\_\_\_} \\
    \textbf{姓名：\_\_\_\_\_\_\_\_\_\_\_\_\_\_\_\_\_\_\_\_\_\_\_\_\_\_\_\_\_\_\_\_\_} \\
    \textbf{学号：\_\_\_\_\_\_\_\_\_\_\_\_\_\_\_\_\_\_\_\_\_\_\_\_\_\_\_\_\_\_\_\_\_} \\[0.5em]
    \textbf{实验日期: \_\_\_\_年\_\_\_\_月\_\_\_\_日\quad 星期\_\_\_上/下午}
\end{center}

\vspace{1.5cm}

\begin{center}
    \fontsize{14pt}{20pt}\selectfont
    浙江大学物理实验教学中心
\end{center}

\newpage

% ========== 正文 ==========

{\fontsize{14pt}{20pt}\selectfont \textbf{1. 实验综述}}

\vspace{0.3cm}

{\fontsize{12pt}{18pt}\selectfont

铁磁材料（如Fe、Co、Ni及其合金）内部存在自发磁化的磁畴结构。在外磁场$H$作用下，磁畴旋转和畴壁迁移使材料被磁化，磁感应强度$B$随$H$增大而增大直至饱和。当$H$周期性变化时，$B$的变化滞后于$H$，形成磁滞回线，其包围面积反映磁滞损耗。磁滞回线的特征参数包括饱和磁感应强度$B_s$、剩余磁感应强度$B_r$和矫顽力$H_c$。

本实验采用\textbf{转换法}，利用示波器的X-Y模式显示B-H曲线。将铁磁样品置于交流励磁电路中，通过电阻$R_1$两端电压反映磁场强度$H$（X通道），通过RC积分电路的电容$C$两端电压反映磁感应强度$B$（Y通道）。逐步改变交流励磁电压，测量不同磁滞回线端点坐标，获得基本磁化曲线；在最大励磁电压下记录完整磁滞回线数据，并用Origin软件绘图和计算磁滞损耗。

}

\vspace{0.5cm}

{\fontsize{14pt}{20pt}\selectfont \textbf{2. 实验重点}}

\vspace{0.3cm}

{\fontsize{12pt}{18pt}\selectfont

掌握磁滞、磁滞回线和磁化曲线的基本概念，理解转换法的原理；学会利用示波器X-Y模式测绘B-H曲线，正确读取饱和磁感应强度$B_s$、剩磁$B_r$和矫顽力$H_c$等特征参数。

}

\vspace{0.5cm}

{\fontsize{14pt}{20pt}\selectfont \textbf{3. 实验难点}}

\vspace{0.3cm}

{\fontsize{12pt}{18pt}\selectfont

正确搭建交流励磁和RC积分电路，理解示波器电压信号到$H$和$B$的量纲转换关系；退磁操作的规范性（电压须单调降至0V）；利用Origin软件计算磁滞回线包围面积以求磁滞损耗。

}

\vfill

% ========== 注意事项 ==========
\noindent\rule{\textwidth}{0.4pt}

{\CJKfontspec[AutoFakeBold]{STFangsong}\fontsize{12pt}{18pt}\selectfont
\textbf{注意事项：}
\begin{enumerate}[label=\arabic*., leftmargin=2em, nosep]
    \item 用PDF格式上传"预习报告"，文件名：学生姓名+学号+实验名称+周次。
    \item "预习报告"必须递交在"学在浙大"的本课程的对应实验项目的"作业"模块内。
    \item "预习报告"还须拷贝到"实验报告"中（便以教师批改）。
\end{enumerate}
}

\vspace{0.5cm}
\begin{center}
    {\CJKfontspec[AutoFakeBold]{STFangsong}\fontsize{12pt}{18pt}\selectfont \textbf{浙江大学物理实验教学中心制}}
\end{center}

\end{document}
